\documentclass{ca-capstone}
\capstonerefs{crandall-pomerance,welschenbach,knuth-taocp2,sagemath}
\begin{document}
% Generated by `make bench CAPSTONE=montgomery-exponentiation` (taskgen.capstone report) from
% bench/results.json: reference time x slack on the reference machine.
% Do not edit: change bench/config.yaml and re-run the benchmark.
\newcommand{\ReqSlack}{10}
\newcommand{\ReqBitsMax}{1024}
\newcommand{\ReqSeconds}{0{,}5}

\capstonetitle{Арифметика Монтгомери и быстрое возведение в степень}

\goal
Требуется реализовать модульное умножение без тривиальных делений и сравнить основные методы
быстрого возведения в степень по модулю.

\functionality
\begin{enumerate}
  \item Редукция Монтгомери (REDC), умножение и возведение в степень для допустимых нечётных
    модулей $N$.
  \item Бинарный алгоритм <<слева-направо>> вычисления $g^x \pmod{N}$.
  \item $q$-арный алгоритм для $q=2^k$ и метод скользящего окна.
  \item Предвычисления степеней фиксированного основания $g$ и их повторное использование между
    запросами; смена основания, модуля, параметров окна или представления чисел не должна приводить
    к использованию неподходящего старого кэша.
  \item Подсчёт возведений в квадрат, умножений и REDC, а также сравнительный эксперимент:
    методы, размеры окна, новые и повторно используемые предвычисления, обычная и
    монтгомериевская модульная арифметика.
  \item Промежуточные трассировки REDC и разложения показателя.
\end{enumerate}

\requirements
\begin{enumerate}
  \item Результаты должны совпадать с библиотечным модульным возведением в степень.
  \item Должно быть определено поведение для $N=1$, $x=0$, $g=0$  и неверных параметров окна. 
  Метод Монтгомери обязан отклонять чётные $N$.
  \item Примеры с модулем и показателем длиной до \textbf{\ReqBitsMax~бит} должны выполняться
    не более чем за \textbf{\ReqSeconds~с} на запрос.
\end{enumerate}

\details
Разрешается использовать любой язык программирования, в частности Python, систему 
компьютерной алгебры SageMath и встроенные большие целые. Трёхаргументный \texttt{pow} и
библиотечные реализации Монтгомери в решении запрещены, но используются как оракул в тестах.
Операции сложения, умножения и деления больших чисел сами не являются предметом работы.

Предвычисление и выполнение основной операции должны измеряться раздельно. Интерфейс должен 
позволять повторно возводить одно основание в разные степени и контролировать попадания в кэш.

\schedule
Работа выполняется индивидуально. Срок выполнения -- \textbf{месяц}; ожидаемая трудоёмкость
-- примерно 36-40 часов. Реализация должна быть разработана и написана вами.

\remarks
\begin{enumerate}
  \item Обязательны unit-тесты REDC и всех методов на случайных и краевых входах.
  \item Следует проверить как повторные обращения к одному основанию, так и смену основания
    после заполнения кэша.
  \item Оптимизация кода не требуется.
\end{enumerate}

\references
\end{document}
